\subsection{Collider Signature of Bulk Neutrinos in Large Extra Dimensions --- arXiv:hep-ph/0405220}

\textbf{Authors:} Qing-Hong Cao, Shrihari Gopalakrishna, C.-P. Yuan

\paragraph{主な主張}
Bulk right-handed neutrino の多数の KK mode を終状態に足し上げることで、特に余剰次元数が $\delta>2$ の場合に、isolated high-$p_T$ charged lepton と missing transverse energy を伴う Higgs production が増強され得て、電子・ミューオン channel の事象数の非対称性から neutrino mass hierarchy と絶対質量スケールを調べられることを示した \cite{Cao:2004BulkCollider}。

\paragraph{新規性}
Bulk-neutrino LED に固有の KK multiplicity enhancement を具体的な hadron-collider signature に結び付け、flavor-dependent event rate を neutrino hierarchy・絶対質量スケールの情報へ接続した点にある。

\paragraph{位置付け}
Bulk neutrino の collider phenomenology を代表する初期研究であり、現在の $T^2$/IceCube 研究との関係では、低い KK shell の離散スペクトルを利用する oscillation probe に対して、多数の KK mode を inclusive に足し上げる collider probe を比較する際の基準文献となる。
